% -*- coding: utf-8 -*-

\chapter{Discussion}

\section{Key Findings}

\subsection{Debate Reduces Majority-Class Collapse}

On Hadoop1, the single-agent baseline collapses to the majority class (\texttt{machine\_down})
under the strict four-class scheme, achieving 50.9\% accuracy that is entirely attributable to
predicting the dominant class while failing on all minority classes. This pattern is a well-known
failure mode of single-pass generation: without a mechanism to enumerate and compare alternatives,
a language model will commit to whichever explanation the most prominent evidence pattern suggests,
which in imbalanced datasets is typically the frequent-class
signature~\cite{du2023improvingfactualityreasoninglanguage}. Macro-\fone{} reveals the severity of
this collapse---the single-agent achieves only 17.3\% macro-\fone{} on the strict scheme, compared
to 21.6\% for the multi-agent system, which correctly identifies \texttt{disk\_full} incidents
(36.4\% per-class \fone{}) while the single-agent achieves 0\%.

This distinction between prediction accuracy and diagnostic utility is practically significant.
Macro-averaged metrics are the appropriate evaluation criterion when rare failure classes are
operationally critical, because they assign equal weight to each class regardless of its
frequency~\cite{le2022logdlhowfar}. The judge mechanism addresses this by requiring each
hypothesis to be scored on evidence quality; a hypothesis that generically attributes failures to
the most common cause will score lower than one that cites specific disk-related error signatures,
forcing the system toward minority-class predictions when the evidence supports them.

\subsection{Cross-Domain Generalization via Structured Verification}

The CMCC results provide the clearest evidence that single-agent generation is insufficient for
cross-domain RCA. Both the single-agent baseline (4.3\% accuracy) and the RAG baseline (8.6\%
accuracy) collapse predominantly to the \texttt{unknown} label---60.2\% and 70.0\% unknown rates
respectively---while the multi-agent system achieves 61.3\% accuracy with only 17.2\% unknown
predictions. The collapse of RAG-based single-agent generation is particularly informative: despite
having retrieved historical context, the single-agent still fails to map its reasoning into the
correct label space. This is consistent with the finding of Roy et al.\ that out-of-distribution
retrieval can mislead rather than ground a single-agent model, because there is no verification
mechanism to assess whether the retrieved evidence is relevant~\cite{roy2024llmagentRCA}.

The multi-agent design addresses this failure mode through structural separation: hypothesis
generation and hypothesis evaluation are handled by different agents operating under different
objectives. The reasoners are tasked with generating diverse candidate explanations, while the
judge is tasked with assessing which candidates are best supported by the available evidence. This
separation is critical in cross-domain settings where the model cannot rely on memorized
correlations from pretraining and must reason explicitly from the evidence provided in the
prompt~\cite{ahmed2023rcarecommend}.

\subsection{HDFS-v1: Balanced Recall in Large-Scale Binary Diagnosis}

HDFS\_v1 tests the system on a large-scale log dataset where the full data distribution is highly
imbalanced (97.1\% normal, 2.9\% anomaly). We evaluate on a balanced 200-sample draw to ensure
meaningful macro metrics across both classes. The single-agent baseline reveals the expected
asymmetry under imbalance: it achieves 97\% recall on normal blocks but only 33\% recall on
anomaly blocks, indicating that it defaults conservatively to the normal prediction when uncertain.
The multi-agent system substantially corrects this asymmetry, achieving 74\% recall on normal and
65\% recall on anomaly, resulting in a balanced macro-\fone{} of 69.4\% versus the single-agent's
61.3\%. This improvement is consistent with the general principle that diverse hypothesis
generation and adversarial verification reduce the pull toward the dominant prediction class,
a mechanism analogous to the ensemble diversity effect documented in the multi-agent debate
literature~\cite{du2023improvingfactualityreasoninglanguage}.

\subsection{Knowledge Graph Grounding Improves Factuality, but Requires Verification}

The comparison between the RAG baseline and the single-agent baseline isolates the contribution of
KG retrieval. On Hadoop1 coarse, retrieval improves macro-\fone{} from 25.8\% to 37.9\%
(+12.1\,pp), confirming that historical incident context provides useful grounding that parametric
model knowledge alone cannot supply. This finding is consistent with prior work showing that RAG
substantially improves factual accuracy for knowledge-intensive tasks by conditioning generation on
retrieved documents~\cite{lewis2021retrievalaugmentedgenerationknowledgeintensivenlp}. However,
on CMCC, where the CMCC label space is more domain-specific, the RAG baseline with 8.6\% accuracy
still dramatically underperforms the multi-agent system at 61.3\%. The retrieved context did not
prevent the single-agent from producing uninformative outputs, because a single-pass model lacks
the verification mechanism to distinguish reliable retrieved evidence from irrelevant
matches~\cite{roy2024llmagentRCA,cui2025aetherlog}. The multi-agent system's judge addresses this
by explicitly scoring the evidence quality criterion, penalizing hypotheses that rely on weakly
supported retrieval.

\subsection{Why Multi-Agent Debate Improves Reliability: A Mechanistic Interpretation}

The consistent improvement of the multi-agent system across all three datasets can be attributed
to three interacting mechanisms. First, specialization across reasoning agents provides
perspective diversity: the log-focused reasoner attends to temporal event patterns and error
signatures in the current incident, while the KG-focused reasoner reasons from historical
recurrence, and the hybrid reasoner synthesizes both. This diversity reduces the risk that all
agents share the same blind spots---a well-documented advantage of ensemble and multi-agent
architectures over single-model predictions~\cite{du2023improvingfactualityreasoninglanguage,wu2023autogen}.
Second, the multi-round debate enables iterative refinement: reasoners receive judge feedback and
the best hypotheses from other agents before generating their next round of candidates. This
cross-pollination allows an agent to incorporate insights it would not have generated
independently, analogous to the self-consistency improvement obtained by aggregating multiple
chain-of-thought traces~\cite{wang2023selfconsistencyimproveschainthought,wei2023chainofthoughtpromptingelicitsreasoning}.
Third, the judge scoring operationalizes reliability as a quantitative selection criterion. Rather
than accepting the first plausible hypothesis, the system enforces a competition in which evidence
quality, reasoning strength, confidence calibration, completeness, and consistency are all
evaluated. This guards against the overconfident but unsupported diagnoses that are a documented
failure mode of single-agent LLM systems~\cite{chen2024rcacopilot}.

\section{Multi-Agent vs.\ Single-Agent Comparison}

\subsection{Quantitative Comparison}
Table~\ref{tab:ma-vs-sa} summarizes the key differences in performance.

\begin{table}[H]
\centering
\caption{Multi-Agent vs.\ Single-Agent Performance Summary}
\label{tab:ma-vs-sa}
\begin{tabular}{lrrr}
\toprule
\textbf{Dataset / Metric} & \textbf{Multi-Agent} & \textbf{Single-Agent} & \textbf{Improvement} \\
\midrule
Hadoop1 Coarse Macro-\fone{} & 39.1\% & 25.8\% & +13.3\,pp \\
Hadoop1 Disk\_full \fone{}   & 36.4\% & 0.0\%  & +36.4\,pp \\
CMCC Accuracy                & 61.3\% & 4.3\%  & +57.0\,pp \\
CMCC Unknown Rate            & 17.2\% & 60.2\% & $-$43.0\,pp \\
HDFS\_v1 Accuracy            & 69.5\% & 65.0\% & +4.5\,pp  \\
HDFS\_v1 Anomaly \fone{}     & 68.1\% & 48.5\% & +19.6\,pp \\
\bottomrule
\end{tabular}
\end{table}

\subsection{Qualitative Differences}

Beyond quantitative metrics, the multi-agent system produces qualitatively different outputs that
better serve operator workflows. Each incident yields 9--15 candidate hypotheses (3--5 per
reasoner), compared to a single hypothesis from the single-agent baseline. This diversity is
not merely cosmetic: the judge's evidence-quality scoring criterion drives the reasoners toward
more specific hypothesis statements that cite concrete log events or retrieved historical facts,
because generic claims score lower~\cite{du2023improvingfactualityreasoninglanguage}. The judge
Score (0--100) also serves as a calibrated confidence proxy that is more informative than the
self-reported confidence values in single-agent outputs, which are known to correlate poorly
with actual correctness in LLM-generated
content~\cite{chen2024rcacopilot}. Finally, the judge feedback narrative---which explains which
evidence was compelling and which was missing---provides transparency into the selection decision
that operators can use to validate the automated diagnosis.

\subsection{When Single-Agent Wins}
The single-agent baseline achieves higher strict accuracy on Hadoop1 (50.9\% vs 21.8\%). This
occurs because the single-agent collapses to the majority class (\texttt{machine\_down}), which
happens to be correct for 50.9\% of cases. However, this ``win'' is misleading: the single-agent
has 0\% recall on three of four classes, making it useless for diagnosing minority failures.

\subsection{Cost-Benefit Tradeoff}

The multi-agent system requires 8--15$\times$ more LLM calls than the single-agent baseline,
translating to a per-incident runtime of 3--5 minutes versus approximately 30 seconds. This
cost is well justified in the scenarios our evaluation targets: cross-domain settings with no
available training data, imbalanced datasets where minority-class failures are operationally
critical, and situations where explanation quality---not just label accuracy---determines the
utility of the automated diagnosis. For well-understood, high-frequency failure modes where a
single-agent achieves high recall (e.g., simple disk-full events with unambiguous error strings),
the additional compute may be unnecessary. A practical deployment strategy would use the judge
score as a confidence signal to decide when to invoke the full debate: low-confidence cases
escalate to multi-agent mode, while high-confidence cases accept the single-agent result. This
adaptive allocation approach is consistent with the cost-aware inference strategies discussed in
the LLM deployment literature~\cite{chen2024rcacopilot}.

\section{Role of Knowledge Graph}

\subsection{RAG vs.\ Single-Agent: Isolating the Retrieval Contribution}

Comparing the RAG baseline to the pure single-agent baseline isolates the marginal contribution
of KG retrieval. On Hadoop1 coarse, retrieval improves macro-\fone{} from 25.8\% to 37.9\%
(+12.1\,pp), a substantial gain that confirms the value of historical incident context for
Hadoop-domain faults where the KG accumulates relevant cases early. On CMCC, the improvement is
smaller in absolute terms (4.3\% to 8.6\%), because the CMCC evaluation begins with an empty KG
and the retrieval index is sparse for the first processed cases. On HDFS\_v1, RAG (63.0\%) is
slightly inferior to the single-agent (65.0\%), suggesting that entity-overlap retrieval does not
consistently identify relevant historical HDFS block traces, and the injected context occasionally
adds noise. These results are consistent with the general finding that RAG is most effective when
the retrieval index contains high-quality, entity-matched historical cases, and least effective in
cold-start or sparse-index conditions~\cite{lewis2021retrievalaugmentedgenerationknowledgeintensivenlp,roy2024llmagentRCA}.

\subsection{Multi-Agent vs.\ RAG: Isolating the Debate and Judging Contribution}

The comparison between the full multi-agent system and the RAG baseline isolates the contribution
of debate and judging on top of retrieval. On Hadoop1 coarse the gain is modest (+1.2\,pp
macro-\fone{}), suggesting that for Hadoop-domain faults, retrieval already supplies most of the
useful context and debate adds incremental benefit at higher compute cost. The picture changes
dramatically for CMCC: multi-agent achieves 61.3\% accuracy versus 8.6\% for RAG (+52.7\,pp).
As noted above, the RAG baseline collapses to ``unknown'' on the majority of CMCC cases despite
having retrieved historical context. This failure mode demonstrates a key limitation of single-pass
generation with retrieval: the model may receive relevant evidence but still fail to integrate it
into a correctly labeled diagnosis without a verification step that checks whether the generated
explanation is actually consistent with the retrieved facts~\cite{roy2024llmagentRCA,cui2025aetherlog}.
For HDFS\_v1, the multi-agent system improves accuracy by 6.5\,pp (69.5\% vs 63.0\%), confirming
that debate and judging provide meaningful gains even for binary classification tasks where the
evidence signal is relatively direct.

\subsection{KG Coverage and Cold-Start Effects}

The effectiveness of KG retrieval is directly tied to the density and relevance of stored
incidents. In our evaluation, the KG is populated incrementally: as the pipeline processes
incidents, high-confidence diagnoses are written back to the graph, enabling later incidents to
benefit from accumulated history. This setup simulates an operational deployment in which the KG
grows over time and becomes more useful as the system processes more incidents. However, it also
means that the first incidents processed encounter an empty or sparse KG and receive no retrieval
benefit---a cold-start problem analogous to the coverage limitations identified in RAG-based
systems operating on new domains~\cite{lewis2021retrievalaugmentedgenerationknowledgeintensivenlp}.
This cold-start effect partially explains the smaller RAG improvement on CMCC compared to Hadoop1:
Hadoop1 cases are processed from a dataset where incident entities share common vocabulary (HDFS
components, job IDs), producing useful entity-overlap matches early, while CMCC cases involve
more diverse OpenStack service identifiers that require more accumulated history before retrieval
becomes effective. Future work should study the KG coverage threshold at which retrieval begins
to provide consistent benefit.

\section{Debate Protocol Effectiveness}

\subsection{Convergence Behavior}

In our experiments, the debate typically reaches a stable hypothesis within two to three rounds.
In round~1, each reasoner generates hypotheses independently based on its designated evidence
source. In round~2, each reasoner receives the judge's feedback on its previous hypotheses along
with the top-scoring hypotheses from the other agents. This cross-pollination typically causes
at least one reasoner to refine its candidate toward a better-supported explanation, raising the
maximum judge score. By round~3, further refinement yields diminishing returns: the maximum
judge score improves by fewer than 5 points in most cases, triggering the stopping criterion.
Cases that fail to converge within three rounds typically involve genuinely ambiguous evidence,
where multiple causal hypotheses are plausible given the available logs and the KG does not
contain a strongly matching historical precedent. These cases are candidates for human escalation
rather than automated resolution, and the low judge scores they produce can serve as a flag for
this purpose---consistent with the confidence-gated escalation strategies proposed for
production LLM-based RCA systems~\cite{chen2024rcacopilot}.

\subsection{Judge Score Distribution and Calibration}

Judge scores span the full range from 0 to 100, but the empirical distribution is concentrated in
two bands. Hypotheses that are ultimately selected as correct predictions typically score between
70 and 90, reflecting strong evidence citation, clear causal reasoning, and appropriate confidence
calibration. Hypotheses that are rejected or correspond to incorrect predictions typically score
between 40 and 60, reflecting partial evidence, vague causal chains, or overconfident assertions
that the judge penalizes under the calibration criterion. Very low scores below 40 are rare and
generally indicate structurally malformed hypotheses (e.g., missing required fields or circular
reasoning). This score distribution suggests that the judge score is a useful confidence proxy:
a threshold around 60 could serve as an operational decision boundary, below which the automated
diagnosis is flagged for manual review. Using structured agent outputs as calibrated confidence
estimates is consistent with emerging approaches to LLM confidence quantification in production
incident management~\cite{chen2024rcacopilot,roy2024llmagentRCA}.

\subsection{Feedback Quality and Operator Utility}

Qualitative inspection of judge outputs across all three datasets reveals consistent patterns that
confirm the rubric is functioning as intended. Feedback for low-scoring hypotheses correctly
identifies missing evidence---for example, flagging that a hypothesis attributes a failure to
network disconnection without citing any specific connectivity error message in the logs. Feedback
for high-scoring hypotheses rewards specificity: citations of named log events, block IDs, or
retrieved historical incidents that share the same failure pattern receive higher evidence quality
scores. Hypotheses that express high confidence in the absence of direct evidence are consistently
penalized under the calibration criterion, which is precisely the behavior needed to guard against
the overconfidence problem identified in single-agent LLM systems~\cite{chen2024rcacopilot}. This
feedback is useful not only for driving refinement in subsequent rounds but also as a standalone
explanation artifact for operator review: an engineer consulting the system output can read the
judge feedback to understand both why the selected hypothesis was preferred and what evidence gaps
remain, enabling a more informed decision about whether to accept the automated diagnosis or
conduct additional investigation.

\subsection{Ablation Considerations}

While we do not conduct formal ablations in this thesis, the results permit qualitative
inferences about the contribution of each component. Removing the judge would eliminate the
verification mechanism and revert the system to selecting an arbitrary hypothesis from the first
reasoner, which the single-agent results suggest would substantially increase majority-class
collapse and reduce minority-class recall. Reducing from three reasoners to one would eliminate
perspective diversity; empirical studies of multi-agent debate show that this diversity is
primarily responsible for the improvement over single-agent generation, since a single agent with
multiple sampling passes provides weaker gains than multiple agents reasoning from different
evidence sources~\cite{du2023improvingfactualityreasoninglanguage}. Removing refinement rounds
would reduce the system to a single-pass debate, which may still outperform the single-agent
baseline through diversity alone but would forego the iterative improvement that feedback enables.
Formal ablations measuring each of these components independently are an important direction for
future work, as they would provide cleaner attribution of performance gains to specific
architectural choices.

\section{Limitations}

Several limitations of the current work should be considered when interpreting the results.
First, label semantics in the Hadoop1 dataset present a construct validity challenge: the
``normal'' label denotes absence of an injected fault rather than a genuinely quiescent system
state. Logs labeled ``normal'' may still contain connectivity messages, heartbeat events, or
resource usage warnings that the reasoners interpret as failure signals, explaining the near-zero
recall on the normal class across all pipelines. This is an inherent property of the dataset
design~\cite{he2020loghub} rather than a system deficiency, but it complicates interpretation
and motivates more careful label semantics in future RCA benchmarks. Second, KG retrieval quality
depends on coverage, which is low at the beginning of evaluation and improves incrementally. This
cold-start problem may understate the benefit of retrieval in a fully populated operational
deployment where the KG has been built over months of incident history. Third, the multi-agent
pipeline requires 8--15 LLM calls per incident compared to 1 for single-agent approaches, a
compute cost that may be prohibitive for strict real-time incident response. Fourth, we do not
train or evaluate supervised ML/DL baselines such as DeepLog or LogBERT, which limits the breadth
of comparison; these methods require dataset-specific training pipelines that are out of scope for
a zero-shot LLM evaluation, but a comprehensive comparison remains an important direction for
future work~\cite{du2017deeplog,guo2021logbert}.

\subsection{Model and prompt sensitivity}
The system relies on prompted generation for both structured extraction and reasoning. Different
foundation models can vary significantly in instruction-following behavior and output formatting,
and even within a model, outputs may be sensitive to prompt wording and context selection. This is
a known characteristic of modern transformer-based
LLMs~\cite{vaswani2023attentionneed,brown2020languagemodelsfewshotlearners}.

While instruction tuning and alignment methods improve usability, they do not eliminate variance.
For example, instruction-following models trained with human feedback improve helpfulness and
adherence to instructions~\cite{ouyang2022traininglanguagemodelsfollow}, and methods such as
Self-Instruct and LIMA explore lightweight data strategies for
alignment~\cite{wang2023selfinstructaligninglanguagemodels,zhou2023limaalignment}. In our
evaluation we mitigate variability through conservative decoding for structured components and by
constraining the number of debate rounds.

\subsection{Cost and latency tradeoffs}
Multi-agent debate increases the number of LLM calls compared to single-agent or RAG baselines.
In practice, this introduces a compute cost that may be acceptable in offline analysis but can be
challenging for strict real-time incident response. A key engineering direction is to adaptively
allocate compute: e.g., run a single-agent baseline for easy cases and escalate to debate only
when uncertainty is high.

\subsection{Knowledge graph construction}
The knowledge graph is an enabling component but also an engineering burden. Building a
high-quality KG requires consistent entity normalization, incident curation, and (ideally)
operator-validated root causes. The benefit of retrieval depends on the presence of relevant
historical incidents; for novel failures or sparse history, retrieval may provide limited value.

\section{Threats to Validity}
We mitigate non-determinism using low temperatures for structured components, but variability
remains. Additionally, normalization from free-text hypotheses to dataset labels introduces
heuristic error.

\subsection{Internal validity (pipeline non-determinism)}
Even with fixed prompts, stochastic sampling can lead to different hypotheses and judge rankings.
This is a general limitation of sampling-based generation. While some work proposes methods for
improving stability via multiple sampled reasoning paths and
selection~\cite{wang2023selfconsistencyimproveschainthought}, large-scale evaluation under such
strategies can be computationally expensive. Our current mitigation is to reduce randomness for
structured steps and to limit the number of debate rounds.

\subsection{Construct validity (label mapping)}
Evaluation requires mapping a generated hypothesis into a discrete label. This mapping is
implemented with conservative rules to avoid over-interpreting ambiguous text, but any rule-based
mapping can introduce errors. This is particularly important for CMCC where symptom-level
hypotheses may correspond to multiple failure classes (e.g., connection failures manifesting as
downstream service errors).

\subsection{External validity (dataset representativeness)}
The datasets cover multiple domains (Hadoop distributed storage, OpenStack cloud management, and
HDFS block traces), but they still represent a subset of real operational environments. Moreover,
benchmark datasets may differ from production in log completeness, noise level, and the quality of
ground-truth labels.

\subsection{Reproducibility and deployment validity}
We emphasize local inference and script-based evaluation for reproducibility. However, model
updates, different quantization settings, and hardware differences can affect latency and output
quality. Parameter-efficient adaptation methods such as LoRA and QLoRA provide promising future
directions for stabilizing domain behavior under limited
resources~\cite{hu2021loralowrankadaptationlarge,dettmers2023qloraefficientfinetuningquantized}.
In addition, preference-optimization approaches such as DPO may offer a simpler alternative to
RLHF pipelines~\cite{rafailov2024directpreferenceoptimizationlanguage}, though applying them
responsibly requires curated preference data.
